% SPDX-License-Identifier: Apache-2.0
\section{An SPI Master, and a Chip to Talk To}
\label{sec:x-spi}

SPI is four wires and no protocol. A select says which chip is
listening, a clock says when a bit is worth looking at, and the two
data wires carry a bit each way on every one of those clocks. What
the bytes mean is between the two ends.

\code{Spi} is three registers. The first holds the divider, the two
mode bits and the chip select; the second is written with the byte to
send and read for the byte that came back; the third says whether a
transfer is running and whether one has finished.

The select is a bit a program sets and clears, rather than something a
transfer does for itself. That is not laziness. Every SPI command is
several bytes under one unbroken select: the read in this example is a
command byte, three of address, one that is thrown away, and then as
many as are wanted. A master that dropped the select between bytes
would start a new command each time and read nothing.

The second register is the one that surprises people, and it is worth
saying why it is one register and not two. SPI has no direction. The
eight clocks that send a byte are the eight that receive one, so
every transfer is a swap: a driver writes a byte it may not care about
in order to clock in a byte it does. Reading this flash means sending
five bytes of command and then sending four more zeroes to get four
bytes back.

\subsection{The Two Mode Bits, and the One Place They Are Not Symmetric}

\code{cpol} is the level the clock idles at. \code{cpha} says which of
a bit's two edges carries it: with \code{cpha} clear the master takes
the bit on the leading edge and moves its own data on the trailing
one, and with \code{cpha} set it does the reverse.

Written that way the two look symmetric, and the implementation is
almost symmetric, but not quite. With \code{cpha} set the first bit
has to be on the wire \emph{before} the first leading edge, and that
leading edge must not move it, or the last bit of the byte never gets
sent. The run counts half edges rather than bits, which makes the
exception a single condition: it is the one leading edge that does not
shift, the one at count zero.

\subsection{What Checks It}

A model of a flash chip, \code{FlashDevice}, on the other end of the
wires, stepped between cycles from the loop as the I2C device of
\code{//docs:hdmi} is. It answers the two commands a bootloader needs:
\code{0x9f}, which says who the chip is, and \code{0x0b}, a fast read.
It is a model and not a part. It holds a \code{Vec} and nothing lowers
it, and its place is a testbench.

It is also the thing that caught the first version of this example.
The identify command came back as \code{de 80 30} against the
\code{ef 40 18} that was asked for, every byte shifted up by one bit,
because the model moved its output on the edge after the one the
master read it on and so lost the top bit of every byte. A slave
presents a bit and moves it only once the master has taken it. Two
shift registers turning on the same clock have to agree about which
edge does which, and the failure when they do not is not an error
anywhere: it is a plausible number that is wrong.

\begin{figure}[htbp]
\centering
\input{spi_timing.tex}
\caption{A window on the identify command: the select held low, the
clock running at a half bit of two cycles, the byte going out on
\code{mosi} and the chip's answer arriving on \code{miso}, a bit per
leading edge.}
\label{fig:x-spi}
\end{figure}

\lstinputlisting[caption={\code{ex\_spi.rs}. Run with \code{bazel run //lib/examples:ex\_spi}.},label={lst:x-spi}]{lib/examples/ex_spi.rs}

\lstinputlisting[style=ascii,caption={What \code{ex\_spi} printed when the build ran it: the chip's identity, the four bytes read from an address, and then the master's Verilog.},label={lst:x-spi-out}]{out_spi.txt}
